% Adaptado em 2026-11-01 para Leonardo Rodrigues.
% Alterações: conteúdo, fonte Unicode 11pt para XeTeX, entradas sem tabelas, resumo e duas páginas.
% Origem: https://github.com/devcelio/resume-template
% Compilar localmente: tectonic cv.tex
% GitHub Repo and Documentation: https://github.com/celiobjunior/resume-template
% Copyright © 2025 Celio B Junior. All rights reserved.
% 
% Licensed under the Apache License, Version 2.0 (the "License");
% you may not use this file except in compliance with the License.
% You may obtain a copy of the License at
%
%     http://www.apache.org/licenses/LICENSE-2.0
%
% This template follows best practices from README.md

% Início do documento LaTeX: define tipo e formato do documento
% a4paper = tamanho A4, 10pt = tamanho base da fonte
\documentclass[a4paper,11pt]{article}

% --- PACOTES: BASE / IDIOMA ---
\usepackage{fontspec}
\setmainfont{lmroman10-regular.otf}[
  BoldFont=lmroman10-bold.otf,
  ItalicFont=lmroman10-italic.otf,
  BoldItalicFont=lmroman10-bolditalic.otf
]
\usepackage[portuguese]{babel}

% --- PACOTES: LAYOUT / TIPOGRAFIA ---
\usepackage{geometry}
\usepackage{parskip}
\usepackage{microtype}
\usepackage{enumitem}
\usepackage{titlesec}
\usepackage{array}
\usepackage{tabularx}

% --- PACOTES: LINKS / BOOKMARKS ---
\usepackage{hyperref}
\usepackage{bookmark}
\usepackage{xurl}

% --- CONFIGURAÇÃO: MARGENS / ESPAÇAMENTO ---
\geometry{top=1.5cm, bottom=1.5cm, left=1.5cm, right=1.5cm} % Margens da página
\setcounter{secnumdepth}{0}                                 % Remove numeração de seção
\setlist[itemize]{
    leftmargin=0.75em, % Recuo da lista (mais perto da margem da página)
    itemsep=0.2em,     % Espaço entre itens
    topsep=0.25em,     % Espaço antes/depois da lista
    parsep=0em,        % Espaço entre parágrafos no mesmo item
    partopsep=0em      % Espaço extra ao iniciar a lista
}

% Remove números de página e cabeçalhos para visual limpo do currículo
\pagestyle{empty}

% --- CONFIGURAÇÃO: METADADOS / LINKS ---
\hypersetup{
    pdftitle={Currículo Leonardo Rodrigues},   % Título do PDF
    pdfauthor={Leonardo Rodrigues},     % Autor do PDF
    colorlinks=true,          % Usa cor nos links (sem caixa)
    linkcolor=black,          % Cor de links internos
    urlcolor=black,           % Cor de URLs
    citecolor=black,          % Cor de citações
    bookmarksdepth=2          % Barra lateral: seção e subseção
}

% --- CONFIGURAÇÃO: TÍTULOS ---
\titleformat{\section}
{\Large\bfseries}
{}
{0em}
{}
[\titlerule\vspace{0.5ex}]
\titleformat{\subsection}
{\normalsize\bfseries}
{}
{0em}
{}
\titlespacing*{\subsection}{0pt}{0.35em}{0.15em}

% Macro para entradas do currículo com bookmark no PDF
\newcounter{cventry}
% Entradas em fluxo textual para preservar a mesma ordem entre leitores de PDF.
\newcommand{\cventry}[4]{%
    \refstepcounter{cventry}%
    \phantomsection%
    \pdfbookmark[2]{#1}{cventry-\thecventry}%
    \noindent\textbf{#1}\par
    \if\relax\detokenize{#3}\relax\else
      \noindent\textit{#3}\par
    \fi
    \if\relax\detokenize{#4}\relax
      \if\relax\detokenize{#2}\relax\else\noindent\textit{#2}\par\fi
    \else
      \noindent\textit{#4}%
      \if\relax\detokenize{#2}\relax\else\textit{ · #2}\fi\par
    \fi
}

\begin{document}
\begin{center}
{\LARGE\textbf{Leonardo Rodrigues}}\\[0.1cm]
\textbf{Desenvolvedor Back-end e Engenheiro de Dados Sênior}\\[0.1cm]

            São Paulo, Brasil ·
            \href{mailto:leo_rodrigues1999@hotmail.com}{leo\_rodrigues1999@hotmail.com}
          \\[0.05cm]

            \href{https://www.linkedin.com/in/leogues}{linkedin.com/in/leogues} ·
            \href{https://github.com/leogues}{github.com/leogues} ·
            \href{https://www.leogues.com}{leogues.com}
          
\end{center}
\section{Resumo profissional}

          Desenvolvedor back-end e engenheiro de dados sênior, com experiência desde 2020.
          Atuação com Go, Node.js/TypeScript e Python em serviços distribuídos e pipelines
          de analytics, com Kafka, ClickHouse e Dagster. Foco em performance,
          confiabilidade e otimização de custos em ambientes GCP e AWS.
        
\section{Experiência profissional}
\cventry{Nemu}{}{Engenheiro de software sênior}{setembro de 2024 – atual}

            Arquitetura e evolução de sistemas distribuídos e plataformas de dados para
            analytics e atribuição de marketing no e-commerce. A plataforma atende
            marcas como Netshoes, WePink, C\&A, Logitech e Nestlé.
          
\begin{itemize}
\item 
              Reduzi em aproximadamente \textbf{90\% o uso de memória do Redis},
              de 50 GB para 5,5 GB, com compressão ZSTD.
            
\item 
              Simplifiquei pipelines com ClickHouse, acelerando a ingestão e o
              processamento, com consultas analíticas
              \textbf{até 60 vezes mais rápidas} em comparação ao MySQL.
            
\item 
              Migrei \textbf{bilhões de registros} do MySQL para o ClickHouse
              com ferramenta interna em Go, leitura paralela por faixas lexicográficas
              e escrita de aproximadamente \textbf{700 mil registros por segundo}.
              Implementei a sincronização contínua de novos dados e alterações via CDC
              e Kafka, utilizando ReplacingMergeTree para consolidar versões dos registros.
            
\item 
              Redução de até 30\% nos gastos totais com otimizações pontuais de FinOps.
            
\item 
              Reduzi em aproximadamente \textbf{95\% o armazenamento analítico},
              de 18,1 TB para 910 GB, com compressão nativa do ClickHouse. A menor demanda
              de IOPS e banda também permitiu adotar discos mais econômicos.
            
\item 
              Desenvolvi e operei pipelines e serviços com Dagster (Python),
              ClickHouse, Kafka e Node.js/TypeScript na GCP, com GKE, MySQL e Memorystore
              (Redis), infraestrutura como código e entrega contínua via GitOps (Argo CD).
            
\item 
              Implantei e operei ClickHouse e Kafka no GKE, com replicação
              multi-AZ, réplicas em nós distintos, backups, alertas e dashboards.
            
\item 
              Investiguei falhas e otimizei a performance de serviços com profiling,
              análise de planos de execução, logs, métricas e traces da stack LGTM
              (Loki, Grafana, Tempo e Mimir/Prometheus).
            
\end{itemize}
\newpage
\section{Experiência profissional — continuação}
\cventry{Engenheiro de software}{São Paulo, Brasil}{}{novembro de 2020 – setembro de 2024}
\begin{itemize}
\item 
              Desenvolvi e aprimorei aplicações back-end em Node.js, com
              melhoria da arquitetura, qualidade de código e facilidade de manutenção.
            
\item 
              Identifiquei gargalos e otimizei a eficiência das aplicações com
              profiling.
            
\item 
              Configurei e operei infraestrutura na AWS: EKS, ECS/Fargate,
              API Gateway, balanceadores de carga, VPC, subnets, rotas e gateways.
            
\item 
              Configurei e mantive PostgreSQL e MySQL no Amazon RDS, além de
              MongoDB. Implementei e aprimorei pipelines de integração e entrega
              contínua (CI/CD).
            
\end{itemize}
\cventry{TM Shop}{São Paulo, Brasil}{Sociedade e-commerce}{fevereiro de 2022 – novembro de 2023}
\begin{itemize}
\item 
              Desenvolvi e mantive integrações com o ERP Bling e criei
              ferramentas internas para apoiar as operações da equipe.
            
\item 
              Levantei requisitos e implementei automações para reduzir
              tarefas manuais, padronizar rotinas e melhorar processos internos.
            
\item 
              Organizei informações de produtos para apoiar a criação de anúncios
              e as operações comerciais.
            
\end{itemize}
\section{Competências técnicas}
\textbf{Back-end:} Go (Golang), Node.js, TypeScript, NestJS e Python.\par

            \textbf{Dados e mensageria:} ClickHouse, Apache Kafka, Dagster,
            PostgreSQL, MySQL, MongoDB, Redis, CDC e modelagem analítica.
          \par

            \textbf{Google Cloud (GCP):} GKE, Cloud SQL (MySQL e PostgreSQL), Memorystore
            (Redis), Cloud Storage, Artifact Registry, IAM, Secret Manager e Cloud KMS.
          \par

            \textbf{GKE e Kubernetes:} clusters regionais, node pools,
            autoscaling, nós Spot, Workload Identity, Docker, serviços,
            Ingress/Gateway API e volumes persistentes.
          \par

            \textbf{Redes na GCP:} VPC, sub-redes, rotas, regras de firewall,
            Cloud NAT, VPC Peering, Cloud DNS e balanceamento de carga.
          \par

            \textbf{AWS:} EKS, ECS/Fargate, RDS (PostgreSQL e MySQL),
            API Gateway, VPC, sub-redes, rotas e balanceadores de carga.
          \par

            \textbf{Infraestrutura como código e entrega:} Terraform/OpenTofu,
            Terragrunt, Helm, Argo CD, GitOps e CI/CD.
          \par

            \textbf{Observabilidade:} Grafana, Loki, Tempo, Mimir, Prometheus,
            profiling, análise de planos de execução e testes de carga.
          \par

            \textbf{Arquitetura e qualidade:} sistemas orientados a eventos,
            idempotência, Clean Architecture, DDD, TDD e testes unitários, de integração
            e de contrato.
          \par
\section{Formação acadêmica}
\cventry{Centro Universitário de Rio Preto}{}{Engenharia da Computação}{}
\end{document}
